\documentclass[11pt,a4paper,parskip=half-]{scrartcl}

%-------------------------
% Packages
%-------------------------
\usepackage[utf8]{inputenc}
\usepackage{lmodern}
\usepackage{microtype}
\usepackage[english,czech]{babel}
\selectlanguage{czech}

\usepackage{amsmath, amssymb, amsfonts, mathtools}
\usepackage{amsthm}

\usepackage{algorithm}
\usepackage{algpseudocode}
% \usepackage{minted}

\usepackage{fnpct}
\usepackage[onehalfspacing]{setspace}
\usepackage{enumitem}

\usepackage{geometry}
\geometry{margin=1in}

\usepackage{authoraftertitle}

\usepackage[breaklinks]{hyperref}
\usepackage{cleveref}

% make footnotes rules as wide as the text
\renewcommand{\footnoterule}{%
  \kern -4pt
  \hrule width \textwidth height 0.4pt
  \kern 3.6pt
}

% define helper command for typesetting inline code
\newcommand{\code}[1]{\texttt{#1}}

\setstretch{1.125}

%-------------------------
% Title info
%-------------------------
\title{Piškvorky}
\author{Pavel Altmann}

\hypersetup{
  pdfborderstyle={/S/U/W 0.5},
  pdftitle={\MyTitle},
  pdfauthor={\MyAuthor}
}
\urlstyle{same}

\begin{document}
\maketitle

\newpage
\tableofcontents
\newpage

\section{Zadání}

Cílem projektu bylo naprogramovat v prostředí systému řízení báze dat Oracle
databázovou aplikaci pro hru Piškvorky pomocí jazyka PL/SQL.

\section{Popis aplikace}

Aplikace nemá vnější uživatelské rozhraní, je dostupná v prostředí systému
řízení báze dat prostřednictvím příkazů DML (Data Manipulation Language).
Základní funkce aplikace jsou:

\begin{itemize}
  \item Vytvořil profilu hráče vložením záznamu do tabulky \code{HRAC}
  \item Vytvořit novou hru vložením záznamu do tabulky \code{HRA}
  \item Zobrazit herní desku čtením z pohledu \code{PAPIR}
  \item Zahrátí tahu vložením záznamu do tabulky \code{TAH}
    \begin{itemize}
      \item Zahrnuje automatickou kontrolu správnosti tahu vzhledem k
        parametrům hry
    \end{itemize}
  \item Zobrazení her, které vyhrál začínající hráč, čtením z pohledu
    \code{VYHRY\_ZACINAJICI}
\end{itemize}

\section{Implementace}

\subsection{Datový model}

Datový model je implementován v souboru \code{model.sql} a skládá se z 4
tabulek: \code{HRAC}, \code{HRA}, \code{TAH}, \code{OMEZENI} a číselníku
\code{STAV}. Všechny datové tabulky používají umělý primární klíč -- automaticky
generovanou posloupnost číselných hodnot -- a mají co nejvíce vlastností
kontrolovaných pomocí integritních omezení.

\paragraph{Tabulka \code{HRAC}.} Obsahuje informace o hráčích registrovaných v
aplikaci -- jméno a statistiky výher, proher a remíz.

\paragraph{Tabulka \code{HRA}.} Obsahuje informace o všech rozehraných i již
dohraných hrách -- hrající hráči, parametry hry (rozměry hrací plochy a
minimální délka výherní řady) a údaje o průběhu hry -- časový záznam a stav hry
(zda je rozehraná, popř. jak dopadla).

\paragraph{Tabulka \code{TAH}.} Obsahuje záznamy o tazích -- ke které hře
patří, na jaké pole byl zahrán, kolikátý tah dané hry to byl, časovou značku a
který hráč jej zahrál.

\paragraph{Tabulka \code{OMEZENI}.} Obsahuje informace o minimálních a
maximálních povolených hodnotách nastavitelných parametrů hry.

\paragraph{Číselník \code{STAV}.} Obsahuje seznam stavů hry -- rozehraná, výhra,
prohra, remíza.

\subsection{Pohledy}

V souboru \code{pohledy.sql} jsou implementovány 4 pohledy využívané v aplikaci:
\code{VYHRY\_ZACINAJICI}, \code{PROHRY\_ZACINAJICI}, \code{REMIZE} a
\code{PAPIR}.

\paragraph{Pohled \code{PAPIR}.} Zobrazuje řádky herních desek v člověkem
čitelném formátu vytvořeném funkcí \code{RADEK\_PAPIRU}.

\paragraph{Pohled \code{VYHRY\_ZACINAJICI}.} Zobrazuje seznam her, které
vyhrál začínající hráč.

\paragraph{Pohled \code{PROHRY\_ZACINAJICI}.} Zobrazuje seznam her, které
prohrál začínající hráč.

\paragraph{Pohled \code{REMIZY}.} Zobrazuje seznam her, které skončily remízou

\subsection{Funkce}

V aplikaci je využívána pouze jedna funkce \code{RADEK\_PAPIRU} definovaná v
souboru \code{funkce.sql}.

\paragraph{Funkce \code{RADEK\_PAPIRU}.} Vrací jeden řádek herního desky
formátovaný jako řetězec znaků \code{X}, \code{O} a mezery.

\subsection{Procedury}

Aplikace využívá pouze jednu proceduru \code{ZABRAN\_TAHU} definovanou v souboru
\code{procedury.sql}.

\paragraph{Procedura \code{ZABRAN\_TAHU}.} Kontroluje správnost tahu a případně
vytvoří výjimečnou situaci s odpovídající chybovou hláškou. Jedná se například o
tah mimo hrací plochu, nebo tah na již obsazené pole.

\subsection{Triggery}

V souboru \code{triggery.sql} je definován jediný použitý trigger
\code{TRG\_TAH\_VALIDACE}.

\paragraph{Trigger \code{TRG\_TAH\_VALIDACE}.} Před vložením záznamu o tahu do
tabulky \code{TAH}, zkontroluje jeho správnost pomocí volání procedury
\code{ZABRAN\_TAHU}.

\section{Scénáře}

Za učelem otestování aplikace byly vytvořeny 3 scénáře. První scénář je
správný průběh hry s výhrou začínajícího hráče. Druhý scénář je správný průběh
hry s výhrou druhého hráče. Tyto scénáře kontrolují základní funkcionalitu
aplikace. Třetí scénář potom testuje veškeré varianty nesprávných tahů a tím
kontroluje fungování procedury \code{ZABRAN\_TAHU}.

Každý z scénářů je implementován v jednom souboru s předponou \code{scenar\_}.

\end{document}
